\section{Introduction}
\label{sec:intro}

Breast cancer remains the most commonly diagnosed malignancy and a leading cause of cancer-related mortality among women worldwide \citep{sung2021global}. Despite significant advances in early detection, surgical techniques, and systemic therapies, the heterogeneity of breast cancer---manifested at the molecular, cellular, and microenvironmental levels---continues to pose fundamental challenges to achieving durable responses and cures \citep{perou2000molecular}. The tumor microenvironment (TME), once regarded as a passive bystander, is now recognised as an active participant in tumor initiation, progression, metastasis, and therapeutic resistance \citep{hanahan2011hallmarks}.

Among the diverse cellular components of the TME, cancer-associated fibroblasts (CAFs) have emerged as a particularly abundant and functionally versatile population \citep{kalluri2016biology}. In breast cancer, CAFs can constitute up to 80\% of the tumor stroma, far exceeding their normal counterparts in healthy mammary tissue \citep{orimo2005stromal}. These activated fibroblasts differ from their quiescent normal counterparts in morphology, gene expression, and functional capabilities, actively promoting tumor progression through multiple mechanisms including extracellular matrix (ECM) remodelling, paracrine signalling, metabolic reprogramming, and immune modulation \citep{sahai2020framework}.

Historically, CAFs were viewed as a relatively homogeneous population of tumor-supportive cells. However, the application of single-cell RNA sequencing (scRNA-seq) and spatial transcriptomics technologies has fundamentally challenged this view, revealing an unexpected degree of heterogeneity within the CAF compartment \citep{bartoschek2018spatially, lambrechts2018phenotype}. Multiple transcriptionally distinct CAF subpopulations---including myofibroblastic CAFs (myCAFs), inflammatory CAFs (iCAFs), and antigen-presenting CAFs (apCAFs)---have been identified across breast cancer subtypes, each with distinct spatial distributions, functional roles, and clinical implications \citep{erve2024}.

This heterogeneity has profound implications for both fundamental biology and clinical practice. Different CAF subpopulations may promote or restrain tumor growth, enhance or suppress immune responses, and confer sensitivity or resistance to specific therapies \citep{chhabra2023}. Understanding the full spectrum of CAF diversity---its origins, spatial organization, functional specialization, and therapeutic vulnerabilities---is therefore essential for developing effective CAF-targeted treatment strategies.

Existing reviews have addressed various aspects of CAF biology \citep{kalluri2016biology, sahai2020framework, lavie2022cancer}, but none has systematically integrated findings from the rapidly accumulating single-cell and spatial omics datasets specific to breast cancer. Here we fill this gap by providing a comprehensive review organized around five interconnected themes:

\begin{enumerate}
  \item \textbf{CAF Origins and Activation}: the cellular sources of CAFs and the signalling pathways governing their activation (Section~\ref{sec:origins}).
  \item \textbf{Single-Cell Heterogeneity}: the transcriptionally and functionally distinct CAF subpopulations revealed by scRNA-seq (Section~\ref{sec:heterogeneity}).
  \item \textbf{Spatial Organization}: how spatial transcriptomics and imaging technologies map CAF subtypes to specific tissue niches (Section~\ref{sec:spatial}).
  \item \textbf{CAF--Immune Crosstalk}: the bidirectional signalling between CAFs and immune cells, and its implications for immunotherapy (Section~\ref{sec:immune}).
  \item \textbf{Therapeutic Targeting}: current and emerging strategies for eliminating or reprogramming CAFs (Section~\ref{sec:therapy}).
\end{enumerate}

We further discuss CAFs as mediators of treatment resistance (Section~\ref{sec:resistance}), their potential as biomarkers (Section~\ref{sec:biomarkers}), and identify critical gaps and future directions (Section~\ref{sec:future}). By integrating evidence across experimental models, clinical datasets, and computational analyses, we aim to provide a unified framework for understanding CAF biology in breast cancer and to guide the development of CAF-targeted therapies.
